\chapter{Charakterisierungen und Existenzfälle}
\label{b13guide:ch:characterizations}

\section{Kandidaten prüfen und Werte erkennen}

\begin{theorem}[Universelle Charakterisierung]
\label{b13guide:thm:universal}
Für $S\subseteq P$ und $m,s\in P$ gilt
\[
\begin{aligned}
 s=\sup S&\ \Longleftrightarrow\
     \forall p\in P\;(s\leq p\ \Longleftrightarrow\ p\in\Hi(S)),\\
 m=\inf S&\ \Longleftrightarrow\
     \forall p\in P\;(p\leq m\ \Longleftrightarrow\ p\in\Lo(S)).
\end{aligned}
\]
Insbesondere, bei Existenz der jeweiligen Werte,
\[
 \sup S\leq p\ \Longleftrightarrow\ \forall x\in S\;(x\leq p),\qquad
 p\leq\inf S\ \Longleftrightarrow\ \forall x\in S\;(p\leq x).
\]
\end{theorem}

\begin{theorem}[Charakterisierung durch die gesamte Schrankenmenge]
\label{b13guide:thm:principal}
Für $S\subseteq P$ und $m,s\in P$ gilt
\[
 s=\sup S\ \Longleftrightarrow\ \Hi(S)=\mathord\uparrow s,
 \qquad
 m=\inf S\ \Longleftrightarrow\ \Lo(S)=\mathord\downarrow m.
\]
Ein Supremum existiert somit genau dann, wenn die obere
Schrankenmenge die Menge aller Elemente oberhalb eines einzelnen
Elements von $P$ ist. Für das Infimum gilt die duale Aussage.
\end{theorem}

\begin{theorem}[Angenommene Infima und Suprema]
\label{b13guide:thm:attained}
Für $S\subseteq P$ und $p\in P$ gilt
\[
\begin{aligned}
 p=\max S&\ \Longleftrightarrow\ p=\sup S\ \land\ p\in S,\\
 p=\min S&\ \Longleftrightarrow\ p=\inf S\ \land\ p\in S.
\end{aligned}
\]
Insbesondere gilt $\inf\{p\}=p=\sup\{p\}$.
Ein Maximum liefert stets ein Supremum, ein Minimum stets ein Infimum.
\end{theorem}

\Needspace{18\baselineskip}
\begin{theorem}[Zeugen gegen falsche Schranken]
\label{b13guide:thm:witnesses}
Existiert $s=\sup S$, so gilt für jedes $p\in P$
\[
 s\nleq p\quad\Longleftrightarrow\quad
       \exists x\in S\;(x\nleq p).
\]
Existiert $m=\inf S$, so gilt
\[
 p\nleq m\quad\Longleftrightarrow\quad
       \exists x\in S\;(p\nleq x).
\]
Insbesondere gibt es zu $p<s$ ein $x\in S$ mit $x\nleq p$,
und zu $m<p$ ein $x\in S$ mit $p\nleq x$.
Eine stärkere Forderung $p<x$ beziehungsweise $x<p$ folgt
in einer beliebigen Halbordnung daraus nicht.

In der anderen Vergleichsrichtung gelten die strikten Kriterien
\[
\begin{aligned}
 s<p&\ \Longleftrightarrow\ \exists u\in\Hi(S)\;(u<p),\\
 p<m&\ \Longleftrightarrow\ \exists\ell\in\Lo(S)\;(p<\ell),
\end{aligned}
\]
sofern der jeweilige Wert $s=\sup S$ beziehungsweise $m=\inf S$
existiert.

\end{theorem}


\begin{definition}[Minimale obere und maximale untere Schranken]
\label{b13guide:def:minimal-bound}
Ein $u\in\Hi(S)$ heißt \emph{minimale obere Schranke} von $S$, wenn
\[
 \forall v\in\Hi(S)\;(v\leq u\ \Rightarrow\ v=u).
\]
Ein $\ell\in\Lo(S)$ heißt \emph{maximale untere Schranke}, wenn
\[
 \forall k\in\Lo(S)\;(\ell\leq k\ \Rightarrow\ k=\ell).
\]
Diese Bedingungen schließen nur eine echte Verbesserung in der
genannten Richtung aus; sie verlangen keinen Vergleich mit
unvergleichbaren Schranken.
\end{definition}

\begin{theorem}[Minimale Schranken und kleinste Schranken unterscheiden]
\label{b13guide:thm:minimal-bound}
Jedes Supremum ist eine minimale obere Schranke; jedes Infimum
ist eine maximale untere Schranke. Die Umkehrungen gelten im
allgemeinen nicht. Auch die Existenz genau einer minimalen oberen
Schranke beziehungsweise genau einer maximalen unteren Schranke
garantiert die entsprechende Grenze nicht.
Ist $u$ eine minimale obere Schranke von $S$, so gilt
\[
 u=\sup S\quad\Longleftrightarrow\quad
 \forall v\in\Hi(S)\;(u\leq v\ \lor\ v\leq u).
\]
Ist $\ell$ eine maximale untere Schranke von $S$, so gilt
\[
 \ell=\inf S\quad\Longleftrightarrow\quad
 \forall k\in\Lo(S)\;(\ell\leq k\ \lor\ k\leq\ell).
\]
Die Vergleichbarkeit wird ausschließlich mit dem jeweiligen
Schrankenkandidaten verlangt, nicht für alle Elemente von $P$.
\end{theorem}

\section{Leere Menge, ganzer Träger und Zusammenfallen}

\begin{theorem}[Infimum und Supremum der leeren Menge]
\label{b13guide:thm:empty-extrema}
Es gilt einschließlich der jeweiligen Existenzbedingungen
\[
 \sup_P\varnothing=\min P,\qquad
 \inf_P\varnothing=\max P.
\]
Das Supremum der leeren Menge existiert genau dann, wenn $P$
ein kleinstes Element besitzt; das Infimum genau dann, wenn
$P$ ein größtes Element besitzt. Für $P=\varnothing$ existiert
keiner der beiden Werte.
\end{theorem}

\begin{theorem}[Infimum und Supremum des ganzen Trägers]
\label{b13guide:thm:carrier}
Die folgenden Bedingungen sind jeweils äquivalent:
\[
\begin{aligned}
 \Hi_P(P)\ne\varnothing
 &\ \Longleftrightarrow\ \max P\text{ existiert}
 \ \Longleftrightarrow\ \sup_P P\text{ existiert},\\
 \Lo_P(P)\ne\varnothing
 &\ \Longleftrightarrow\ \min P\text{ existiert}
 \ \Longleftrightarrow\ \inf_P P\text{ existiert}.
\end{aligned}
\]
Im jeweiligen Existenzfall gilt
\[
\begin{aligned}
 \sup_P P&=\max P=\inf_P\varnothing,&
 \Hi_P(P)&=\{\max P\},\\
 \inf_P P&=\min P=\sup_P\varnothing,&
 \Lo_P(P)&=\{\min P\}.
\end{aligned}
\]
\end{theorem}

\begin{theorem}[Infimum unter Supremum; Gleichheitsfall]
\label{b13guide:thm:collapse}
Sei $S\ne\varnothing$, und seien $m=\inf S$ und $s=\sup S$
vorhanden. Dann gilt
\[
 \forall x\in S\;(m\leq x\leq s),\qquad m\leq s.
\]
Ferner sind äquivalent:
\[
 m=s,\qquad s\leq m,\qquad S\text{ ist einelementig}.
\]
Im Gleichheitsfall ist $S=\{m\}=\{s\}$.
Enthält $S$ zwei verschiedene Elemente, so gilt $m<s$.
\end{theorem}

\begin{theorem}[Grenzen abgeschlossener Intervalle und einseitiger Bereiche]
\label{b13guide:thm:interval-values}
Für $a,b\in P$ mit $a\leq b$ gilt
\[
\begin{aligned}
 \Lo([a,b]_P)&=\mathord\downarrow a,&
 \Hi([a,b]_P)&=\mathord\uparrow b,\\
 \inf_P[a,b]_P&=a,& \sup_P[a,b]_P&=b.
\end{aligned}
\]
Für jedes $p\in P$ gilt außerdem
\[
 \Hi(\mathord\downarrow p)=\mathord\uparrow p,\qquad
 \Lo(\mathord\uparrow p)=\mathord\downarrow p,
\]
und die folgenden Grenzen existieren:
\[
 \sup_P(\mathord\downarrow p)=p,\qquad
 \inf_P(\mathord\uparrow p)=p.
\]
\end{theorem}
